% Research-only candidate, 2026-11-16.
% NOT included in main.tex. The two finite propositions have paper proofs;
% the transport/rebound hypothesis is OPEN for the actual logarithmic law.
\subsection{Sixteen-residue escape mechanism and a conditional return bound}
\label{subsec:mod16-escape-candidate}
For odd $c\ge-1$, put $T(n)=T_{3,c}(n)$. Fix a shell
$I_L=(L/2,L]\cap\mathbb N$, and distinguish a \emph{local} lower
threshold $L/4$ from the possibly much smaller \emph{global}
first-passage threshold $a_L$. Write
\[
 G_c=\{y\in\mathbb N:y\equiv0\pmod4\text{ or }
                 3y+c\equiv0\pmod{16}\}.
\]
Modulo $16$, $G_c$ comprises the five classes
$0,4,8,12,5c$. In particular, the uniform distribution on all
$16$ classes gives it mass $5/16$; no analogous claim for a
\emph{conditioned entrance distribution} is being made.

\begin{proposition}[Deterministic local escape]
\label{prop:mod16-local-escape-candidate}
Let $c\ge-1$ be odd and $L\ge\max\{16,2|c|\}$. If
$y\in I_L\cap G_c$, then the orbit starting at $y$ reaches
$(0,L/4]$ in at most four steps, and before that first passage
visits $I_L$ at most twice, counting the initial state.
\end{proposition}
\begin{proof}
If $4\mid y$ then $T(y)=y/2\le L/2$ and
$T^2(y)=y/4\le L/4$. Otherwise $16\mid(3y+c)$, so $y$ is odd.
Put $z=(3y+c)/4$, a positive multiple of $4$. The trajectory
begins $y\mapsto2z\mapsto z\mapsto z/2\mapsto z/4$.
Since $y\le L$ and $c\le L$, we have $z\le L$, hence
$z/4\le L/4$. If $2z\le L$, then $2z$ is the only possible
extra visit to the shell, because $z\le L/2$. If $2z>L$,
then $2z$ lies above the shell and $z$ is its only possible
extra visit. No independence hypothesis is used.
\end{proof}

\begin{proposition}[Finite-modulus conditional return criterion]
\label{prop:mod16-transport-candidate}
Fix a probability or subprobability law on stopped trajectories
that start above a global cutoff $a_L\le L/4$. Let $N$ be
its number of visits to $I_L$ before this global first passage,
with $Y_k$ the $k$-th visited shell state on $\{N\ge k\}$.
Let $B_k$ denote a \emph{rebound}: after the first subsequent
passage below $L/4$, the trajectory revisits $I_L$ before
its global first passage. Suppose constants $g>0$ and
$\vartheta<1$ satisfy, for all $k$ with $\Pr(N\ge k)>0$,
\[
 \Pr(Y_k\in G_c\mid N\ge k)\ge g,
 \qquad
 \Pr(B_k\mid Y_k\in G_c,N\ge k)\le\vartheta.
\]
Then with $\epsilon=g(1-\vartheta)>0$,
\[
 \Pr(N\ge k+2)\le(1-\epsilon)\Pr(N\ge k),
 \qquad
 \mathbb E N^2\le {4(2-\epsilon)\over\epsilon^2}.
\]
It suffices for the first inequality to have
$d_{\rm TV}(\mathcal L(Y_k\bmod16\mid N\ge k),U_{16})\le\delta<5/16$,
with $g=5/16-\delta$.
\end{proposition}
\begin{proof}
On $\{N\ge k,Y_k\in G_c,\neg B_k\}$, the preceding
proposition allows at most one more shell visit, so $N<k+2$.
The conditional hypotheses imply that this event has conditional
probability at least $\epsilon$. Thus $p_{k+2}\le(1-\epsilon)p_k$
for $p_k=\Pr(N\ge k)$. Since
$\mathbb E N^2=\sum_{k\ge1}(2k-1)p_k$ and
$p_{2j+1},p_{2j+2}\le(1-\epsilon)^j$, the displayed moment bound
follows from summing the geometric series. The total-variation
claim follows because $G_c$ occupies five residue classes.
\end{proof}

\paragraph{The missing transport estimates.}
For the actual logarithmic first-entrance law, neither the
conditional modulo-$16$ discrepancy bound nor the uniform
\emph{post-local-exit rebound bound} has been proved. In the
first-passage argument of the paper, the global cutoff may lie
far below $L/4$; therefore passing below $L/4$ does \emph{not}
by itself end the counted trajectory. Consequently this
proposition is a proved \emph{conditional criterion}, not a
proof of $T_X=o(X)$.

\begin{proposition}[Obstruction to monotone position-only certificates]
\label{prop:position-monotone-obstruction-candidate}
Fix odd $c\ge-1$. Let $L$ and $a$ be positive integers satisfying
$L\ge16(|c|+1)$ and $64\le a\le L/4$.
In the terminating first-return graph of $I_L=(L/2,L]$
with global cutoff $a$, there exist return edges
$y_+\to z_+$ and $y_-\to z_-$ such that
$y_+<z_+$ and $y_->z_-$.
In particular no potential of the form $U(y)=f_L(y/L)$, with
$f_L$ monotone (in either direction), can obey
$U\ge1+P_IU$ on the whole terminating graph.
\end{proposition}
\begin{proof}
Choose a power of two $M\ge(3L+|c|)/a$ that is minimal
with this property. Then $M\ge16$ and
$M<2(3L+|c|)/a<L/8$. Since $3$ is invertible modulo
$M$, the congruence $3y+c\equiv0\pmod M$ has exactly one
class modulo $M$. The intervals
\[
 J_+=\bigl(L/2,(2L-c)/3\bigr],\qquad
 J_-=\bigl((2L-c)/3,L\bigr]
\]
have lengths greater than $L/8$ under the stated assumptions;
therefore they each contain a representative of this class,
say $y_+$ and $y_-$. Both are odd.
For $y_+$, the next iterate $z_+=(3y_++c)/2$ belongs to
$I_L$ and exceeds $y_+$. For $y_-$, the first iterate
$(3y_-+c)/2$ exceeds $L$ and is even, while the second iterate
$z_-=(3y_-+c)/4$ belongs to $I_L$ and is smaller than $y_-$.
The prescribed divisibility by $M$ ensures that, after the
respective return, repeated halving reaches a value at most
$(3L+|c|)/M\le a$, so these edges belong to the \emph{terminating}
return graph even without any Collatz convergence assumption.
A nondecreasing $f_L$ contradicts the first strict-return
inequality, while a nonincreasing $f_L$ contradicts the second.
\end{proof}

\paragraph{Status.}
The deterministic escape, conditional return criterion, and
monotone-obstruction proofs are supplied here as a research draft.
The uniform conditional hypotheses are \emph{open}. The arithmetic
construction has independent exact-integer regression tests, but the
new claims are not yet kernel-certified in Lean.
